% vim: tw=78 ai spell spelllang=en

\documentclass[10pt]{llncs}
\usepackage{color}
\usepackage{url}
\usepackage{graphicx}
\usepackage{listings}
\usepackage{hyperref}
\usepackage{xspace}
\pagestyle{plain}

\lstset{basicstyle=\ttfamily\scriptsize,escapeinside={(*@}{@*)}}

%{{{ (La)TeX definitions

\newcommand{\nb}[1]{\marginpar{\scriptsize #1}}
\renewcommand\note[1]{\nb{\textcolor{blue}{#1}}}
\newcommand\todo[1]{\nb{\textcolor{red}{#1}}}
\newcommand\tool{\textsc{Symbiotic}\xspace}
\newcommand\klee{\textsc{Klee}\xspace}

%\renewcommand\note[1]{}
%\renewcommand\todo[1]{}

%}}}

\begin{document}

%{{{ Title + Author + Abstract

\title{\tool~2: More Precise Slicing}
% \thanks{This work has been supported by The Czech Science
%     Foundation (GA\v{C}R), grant No.~P202/12/G061.}}
\subtitle{(Competition Contribution)}
\author{Jiri~Slaby \and Jan~Strej\v{c}ek}
\institute{Faculty of Informatics, Masaryk University\\
  Botanick\'a 68a, 60200 Brno, Czech Republic\\
  \email{\{slaby,strejcek\}@fi.muni.cz}}

\maketitle

\begin{abstract}
  \tool~2 keeps the concept and the structure of the original bug-finding
  tool \tool, but it uses a more precise slicing based on a field-sensitive
  pointer analysis instead of field-insensitive analysis of the original
  tool. The paper discusses this improvement and its consequences.  We also
  briefly recall basic principles of the tool, its strong and weak points,
  installation, and running instructions. Finally, we comment the results
  achieved by \tool 2 in the competition.
\end{abstract}

%}}}

%{{{ Verification approach

\section{Verification Approach and Software Architecture}
\label{sec:approach}

Both \tool~\cite{SST13} and \tool 2 implement our verification approach
proposed in~\cite{sse}. The approach combines three standard techniques,
namely \emph{code instrumentation}, \emph{program slicing}~\cite{weiser},
and \emph{symbolic execution}~\cite{king}. While the approach is originally
designed for detection of bugs described by state-machines, \tool 2 so far
support only one kind of bugs: reachability of an \texttt{ERROR}
label. Hence, we briefly recall the approach restricted just to this simple
kind of errors. We explain the structure of the tool simultaneously.
\begin{enumerate}
\item Code instrumentation inserts \texttt{assert(0)} to each \texttt{ERROR}
  label. It is realized by a \texttt{bash} script calling \texttt{sed}. The
  instrumented code is translated to the \textsc{Llvm} bitcode by
  the \textsc{Clang} compiler.
\item Program slicing removes instructions of the instrumented code that do
  not affect reachability of the inserted \texttt{assert(0)}
  statements. This code size reduction is crucial for the overall efficiency
  of \tool 2. The slicer is implemented in C++ as a plug-in for the
  \textsc{Llvm} optimizer \texttt{opt}.
\item Symbolic execution either reaches \texttt{assert(0)}, or correctly
  finishes the execution without reaching \texttt{assert(0)}, or it runs out
  of time or memory etc. These possibilites correspond to answers
  \texttt{FALSE}, \texttt{TRUE}, \texttt{UNKNOWN}, respectively.
  We use symbolic executor \klee~\cite{klee} whose outputs are translated
  to \texttt{TRUE/FALSE/UNKNOWN} by a simple \texttt{bash} script.
\end{enumerate}
The whole pipeline is executed stepwise by another \texttt{bash} script.

One can easily see that the only original and non-trivial part of the tool
is the slicer. And the slicer also contains all improvements of \tool 2
compared to the original version of \tool competing in SV-COMP 2013. We have
fixed several bugs in the original slicer (invalid treatment of several
instructions and functions with variable number of arguments). We have also
substantially improved the precision of the employed pointer analysis, which
is a crucial part of slicing. \todo{dat tam nejaky odkazy na pointer
  analysis implementovanou v \tool 2? Pripadne nejakou odkaz na
  field-independent pointer analysis.}\note{Je to klasicka Andersen+ta
  vylepseni. Nevim, jestli je to nekde dokumentovane explicitne. Kdyz jsem to
  hledal do dizertace, tak jsem nic rozumneho nenasel.} In particular, we have
switched from a field-insensitive to a field-sensitive pointer analysis.
This means that every field of a \texttt{struct} is now handled as a different
pointer target. Similarly, we handle the first 64 elements of each array as
distinct targets. Finally, we have added some type-filters that speed up the
pointer analysis and make its results more accurate. All these improvements
significantly reduce the number of incorrect answers produced by \tool 2. For
example, \tool 2 produces only correct answers for category ProductLines while
\tool produces 131 incorrect answers on the same category.
 
%}}} {{{ Software architecture - Zakomentov�no

% \section{Software Architecture}
% \paragraph{Libraries/External Tools} For the code translation and for
% operations with the \textsc{Llvm} bitcode we use \textsc{Llvm}/\textsc{Clang}
% 3.2\footnote{\url{http://llvm.org}}.  The symbolic executor is
% \klee~\cite{klee} which itself uses the \textsc{Stp} constraint
% solver~\cite{stp}. We obtained both \klee and \textsc{Stp} from the
% respective GIT repositories and used the snapshots.

% \paragraph{Software Structure and Architecture} The architecture described
% in Section~\ref{sec:approach} is summarized in Figure~\ref{fig:pipeline}.
% The slicer is written as a plug-in for the optimizer \texttt{opt} from the
% \textsc{Llvm} suite. It is publicly available in a separate repository at
% \url{https://github.com/jirislaby/LLVMSlicer/}.
% % The slicer improves the symbolic execution considerably.
% For ease of use, all parts of the tool pipeline are one by one run by a
% single script \texttt{runme}.

% \begin{figure}[tb]
% \begin{center}
% \input{pipeline2.pspdftex}
% \end{center}
%   \caption{Pipeline of the tool}
% \label{fig:pipeline}
% \end{figure}

% \paragraph{Implementation Technology} The instrumentation is performed by a
% \texttt{bash} script using \texttt{sed}. The final filter also uses
% \texttt{bash} with the help of \texttt{grep}. The rest of the toolchain is
% written in C++ and compiled using \texttt{gcc}.

%}}}
%{{{ Discussion

\section{Discussion of Strengths and Weaknesses of the Approach}

\textbf{Prepsat, zpresnit, zlepsit. Napsat, ze concurrency a jeste cosi jde
  mimo nas. Napsat neco o Future work? Podivat se jeste do recenzi.}


The strength of the tool lies in its high precision of the answers. In
theory, the only source of potential incorrect answers is the slicer: it can
completely remove an infinite loop %in some cases 
and an unreachable
\texttt{ERROR} label located below the loop may become reachable. However,
there is no such case in the competition benchmarks.  All incorrect answers
%produced by our tool in the competition
are due to imperfection of our implementation.

%In general, 
The main weakness of the tool is a high percentage of
\texttt{UNKNOWN} results coming mainly from high computation
cost of the symbolic execution of programs with loops or recursion. This
problem is relieved by slicing, but there are still many cases where the
sliced code remains complex and the symbolic execution runs out of time or
memory.  Other sources of \texttt{UNKNOWN} results are internal errors of
\klee and general limitations of constraint solving.  We plan to
reduce this weakness by application of the compact symbolic
execution~\cite{CSE} instead of the classic one.

%}}}

\newpage

%{{{ Tool setup and configuration

\section{Tool Setup and Configuration}
Before using \tool 2, ensure that the target system contains 32-bit
libraries and \texttt{llvm} with \texttt{clang}. \texttt{llvm} and
\texttt{clang} have to be in version 3.2 exactly.  Then, \tool 2 can be
downloaded from %:
%\begin{center}
\url{http://sf.net/projects/symbiotic/}.
%\end{center}
As \klee requires to be run from a pre-defined path, we are obliged to
change the current directory to \texttt{/opt/} and \emph{untar} the
downloaded \tool 2 archive there.  Running of the tool is then
straightforward. When the current directory is \texttt{/opt/symbiotic/}, the
tool can be invoked for each \texttt{<benchmark.c>} from the set by
\texttt{./runme <benchmark.c>}. For benchmarks intended for 64-bit, set
\texttt{MFLAG=-m64} environment variable.  The answers provided by \tool are
as required by the competition rules: \texttt{TRUE/FALSE/UNKWNOWN}. If the
result for \texttt{<benchmark.c>} is \texttt{FALSE}, discovered error paths
can be found in \texttt{<benchmark.c>-klee-out/}.

%}}}
%{{{ Software project and contributors

\section{Software Project and Contributors}
\tool 2 was contributed prevalently by the authors of this paper and Marek
Trt\'{\i}k. Jiri Slaby is a contact person. The tool is licensed under the
GNU GPLv2 License.

% We would like to thank our advisors prof. Anton\'in Ku\v{c}era and dr. Jan
% Strej\v{c}ek. Our work is supported by grant no.

%}}}

\bibliographystyle{plain}
\bibliography{sse2014}

\end{document}
